\documentclass[11pt,a4paper]{article}
\usepackage[T1]{fontenc}
\usepackage{lmodern,microtype}
\usepackage{amsmath,amssymb,amsthm,mathtools}
\usepackage[margin=26mm,headheight=14pt]{geometry}
\usepackage{booktabs,array,tabularx,enumitem}
\usepackage{xcolor,fancyhdr,listings,hyperref}
\definecolor{ink}{HTML}{17324D}
\definecolor{accent}{HTML}{176B78}
\definecolor{codebg}{HTML}{F5F7F9}
\hypersetup{colorlinks=true,linkcolor=ink,citecolor=accent,urlcolor=accent,
  pdftitle={A constructive proof of the quartic alpha+n conjecture},
  pdfsubject={Exact realization of every monic integer quartic by a chromatic factor}}
\newtheorem{theorem}{Theorem}[section]
\newtheorem{lemma}[theorem]{Lemma}
\newtheorem{proposition}[theorem]{Proposition}
\newtheorem{corollary}[theorem]{Corollary}
\theoremstyle{definition}
\newtheorem{definition}[theorem]{Definition}
\newtheorem{remark}[theorem]{Remark}
\newcommand{\Z}{\mathbb Z}
\newcommand{\Q}{\mathbb Q}
\newcommand{\N}{\mathbb N}
\newcommand{\pairs}{\mathcal P}
\newcommand{\parts}{\mathcal M}
\newcommand{\fall}[2]{(#1)_{#2}}
\newcommand{\bit}[1]{[#1]_2}
\newcommand{\htf}{\mathcal H}
\DeclareMathOperator{\Gal}{Gal}
\pagestyle{fancy}
\fancyhf{}
\fancyhead[L]{\small\color{ink}Quartic chromatic factors}
\fancyhead[R]{\small 28 November 2026}
\fancyfoot[C]{\small\thepage}
\setlength{\emergencystretch}{3em}
\setlist{nosep,leftmargin=*}
\raggedbottom
\lstset{language=Python,basicstyle=\ttfamily\scriptsize,
  keywordstyle=\color{ink}\bfseries,commentstyle=\color{accent},
  backgroundcolor=\color{codebg},frame=single,rulecolor=\color{black!12},
  breaklines=true,columns=fullflexible,keepspaces=true,
  showstringspaces=false,tabsize=4,aboveskip=10pt,belowskip=10pt}
\title{\color{ink}\bfseries A constructive proof of the quartic~\texorpdfstring{\boldmath $\alpha+n$}{alpha+n} conjecture}
\author{\normalsize Research manuscript}
\date{28 November 2026}

\begin{document}
\maketitle
\begin{abstract}
For every monic quartic $f\in\Z[x]$, we construct positive integers $N,n$
and a finite simple graph $G$, whose vertices are the union of an $N$-clique
and a four-clique, such that
\[
 P_G(q)=\fall qN f(q-n).
\]
Consequently every quartic algebraic integer becomes a chromatic root after
addition of a positive integer. The construction uses a rational quadratic
profile for six palette intersections. Its degree-four and degree-three
contributions to the remaining constant-term equation cancel identically.
All fifteen palette atoms retain strictly positive quadratic margins;
parity choices and a bounded congruence correction then give exact integers.
For $f(x)=x^4+ax^3+bx^2+cx+d$ and
$\htf=10+|a|+|b|+|c|+|d|$, the construction gives
$|V(G)|\le 54\,080\,000\,\htf^4$.
The proof is elementary and does not depend on computational search.
A complete integer-arithmetic implementation and an independent matching
verification are included.
\end{abstract}

\section{The named problem and the result}

Write $P_G(q)$ for the chromatic polynomial of a finite simple graph $G$.
A chromatic root is a complex zero of some $P_G$. Cameron and Morgan's
$\alpha+n$ conjecture asks whether every algebraic integer $\alpha$ admits
a natural number $n$ for which $\alpha+n$ is a chromatic root
\cite[Conjecture 4]{CM}. The problem arose in the 2008 Isaac Newton
Institute programme. The study group established the quadratic case,
and Bohn established the cubic case \cite{CM,Bohn}. This paper addresses the degree-four case of that same
statement, without changing the number, the permitted translation, or the
class of finite simple graphs.

For $N\ge0$, let $\fall qN=q(q-1)\cdots(q-N+1)$, with
$\fall q0=1$. We prove a polynomial realization statement stronger than
the assertion about one root.

\begin{theorem}[Exact quartic realization]\label{thm:main}
Let $f(x)=x^4+ax^3+bx^2+cx+d$ with $a,b,c,d\in\Z$. There exist
positive integers $N,n$ and a connected finite simple graph $G$ on $N+4$
vertices such that
\begin{equation}\label{eq:main}
 P_G(q)=\fall qN f(q-n).
\end{equation}
The graph consists of two cliques of orders $N$ and $4$, with specified
edges between them, and has chromatic number $N$. No irreducibility
assumption on $f$ is needed.

More explicitly, with
\begin{equation}\label{eq:height}
 \htf=10+|a|+|b|+|c|+|d|,\qquad t=5200\htf^2,
\end{equation}
the formulas below produce such a graph and satisfy
\begin{equation}\label{eq:size}
 0<n<t^2,\qquad N+4\le2t^2=54\,080\,000\,\htf^4.
\end{equation}
\end{theorem}

The graph class just described is called a $(4,N)$-\emph{biclique} in
\cite{Bohn,CM,Gilmore}. In that terminology a biclique is the complement
of a bipartite graph; the word does not mean a complete bipartite graph.

\begin{corollary}\label{cor:root}
Every algebraic integer of degree at most four has a positive integer
translate that is a chromatic root.
\end{corollary}
\begin{proof}
If $\alpha$ has monic minimal polynomial $g$ of degree $e\le4$, apply
Theorem~\ref{thm:main} to $f(x)=x^{4-e}g(x)$. Then
$P_G(\alpha+n)=0$.
\end{proof}

\subsection{What is inherited and what is proved here}

The biclique formulas and the use of nonnegative Venn-region sizes are
established tools. Gilmore's cubic construction already uses positive
quadratic leading terms and congruence conditions; its discussion of the
quartic case ends with a system of fourteen inequalities in eleven
variables \cite[Section 6.4, pp.~59--62]{Gilmore}. The available abstract
of Delbourgo and Morgan's algorithm explicitly conditions its output on
existence of the required biclique \cite{DM19}. Their later paper describes
the quartic results as computational evidence \cite[Section 1.1]{DM22}.

The substantive step here is an explicit quartic solution of the integer
realization problem: a rational profile with simultaneous leading-term
cancellation, positive margins for every palette atom, and an exact
integer lift. Neither a change of coordinates nor a numerical existence
test supplies that step. Sections~\ref{sec:construction}--\ref{sec:positive}
prove it for all four integer coefficients. Section~\ref{sec:effective}
makes the size choice completely explicit.

The supplied VMRC manuscript~245 uses a related proof-design principle:
prescribe upper coefficients while keeping lower correction costs smaller
than a positive realization margin \cite[Sections 2--3]{VMRC}. That
independence-polynomial theorem does not imply a chromatic realization
theorem. Here its role is methodological inspiration; the required
chromatic counting and positivity arguments are proved below.

\section{Four palettes and an exact counting formula}\label{sec:palettes}

Let $C$ be an $N$-clique, and let $v_1,v_2,v_3,v_4$ form another clique.
For each $i$, specify a set $A_i\subseteq C$ of vertices
\emph{not} adjacent to $v_i$. All other cross-pairs are edges. After the
vertices of $C$ receive distinct colors, $A_i$ is the set of old colors
available to $v_i$. If $y=q-N$, all four vertices also have the same
$y$ new colors available.

For a nonempty $J\subseteq[4]=\{1,2,3,4\}$, let $m_J$ be the
number of vertices of $C$ belonging to exactly the sets $A_i$ with
$i\in J$. We use no vertices belonging to none of the $A_i$.
Any fifteen nonnegative integers $m_J$ define actual sets and an actual
finite graph in this way.

Introduce integer parameters $L,z_i,D_{ij},T_i,U$ by prescribing
\begin{align}
 |A_i|&=L+z_i, & |A_i\cap A_j|&=L-D_{ij},\label{eq:intersections1}\\
 \left|\bigcap_{j\ne i}A_j\right|&=L-T_i,
 & |A_1\cap A_2\cap A_3\cap A_4|&=L-U.\label{eq:intersections2}
\end{align}
The parameter $T_i$ always indexes the triple \emph{omitting} $i$.

Let $\pairs$ be the six two-element subsets of $[4]$, and let $\parts$
be the three partitions of $[4]$ into two unordered pairs. When
$\{i,j\}\in\pairs$, write $\{k,\ell\}=[4]\setminus\{i,j\}$.

\begin{lemma}[Palette inversion]\label{lem:venn}
The prescribed intersections are realized by the following atom counts:
\begin{align}
 m_{1234}&=L-U,\label{eq:atom4}\\
 m_{[4]\setminus\{i\}}&=U-T_i,\label{eq:atom3}\\
 m_{ij}&=T_k+T_\ell-D_{ij}-U,\label{eq:atom2}\\
 m_i&=z_i+\sum_{j\ne i}D_{ij}-\sum_{j\ne i}T_j+U.\label{eq:atom1}
\end{align}
Thus nonnegativity of these fifteen integers is sufficient, with no
additional set-realizability condition. Their sum is
\begin{equation}\label{eq:N}
 N=L+\sum_i z_i+\sum_{i<j}D_{ij}-\sum_iT_i+U.
\end{equation}
\end{lemma}
\begin{proof}
Create $m_J$ distinct elements of type $J$ and put an element into $A_i$
exactly when its type contains $i$. Equations
\eqref{eq:atom4}--\eqref{eq:atom1} are inclusion--exclusion for the
specified intersections. For example, the intersection $A_i\cap A_j$
has the pair atom, the two triple atoms omitting $k$ and $\ell$, and
the quadruple atom; their sum is $L-D_{ij}$. Summing the atoms containing
a fixed $i$ gives $L+z_i$. The other identities follow in the same way.
Inclusion--exclusion for the union gives \eqref{eq:N}.
\end{proof}

\begin{lemma}[The centered quartic]\label{lem:quartic}
Put $x=y+L=q-N+L$. The number of extensions of a fixed coloring of $C$
is the following polynomial in $x$:
\begin{align}
 F(x)={}&\prod_{i=1}^4(x+z_i)
 -\sum_{\{i,j\}\in\pairs}(x-D_{ij})(x+z_k)(x+z_\ell)\notag\\
 &+\sum_{\{\{i,j\},\{k,\ell\}\}\in\parts}(x-D_{ij})(x-D_{k\ell})
 +2\sum_i(x-T_i)(x+z_i)-6(x-U).\label{eq:F}
\end{align}
Consequently
\begin{equation}\label{eq:Pfactor}
 P_G(q)=\fall qN F(q-N+L).
\end{equation}
\end{lemma}
\begin{proof}
The common palette of a set of small-clique vertices has size $x+z_i$
for a singleton, $x-D_{ij}$ for a pair, $x-T_i$ for a triple omitting
$i$, and $x-U$ for all four. Count assignments in which all four colors
are distinct by inclusion--exclusion on the partition lattice. A block of
size $r$ contributes the factor $(-1)^{r-1}(r-1)!$. The five partition
types $1+1+1+1$, $2+1+1$, $2+2$, $3+1$, and $4$ give precisely
\eqref{eq:F}. This argument is a specialization of the standard biclique
formula \cite[Proposition 9]{CM}.
It applies to every integer $q\ge N$; hence it is a polynomial identity.
Multiplication by the $\fall qN$ colorings of $C$ proves
\eqref{eq:Pfactor}.
\end{proof}

Let $S=\sum_i z_i$, and let $e_j=e_j(z_1,z_2,z_3,z_4)$ denote the
elementary symmetric functions. Define
\begin{equation}\label{eq:M}
 M=e_4+\sum_{\{i,j\}\in\pairs}D_{ij}z_kz_\ell
     +D_{12}D_{34}+D_{13}D_{24}+D_{14}D_{23}.
\end{equation}
Expansion of \eqref{eq:F} gives
\begin{align}
 [x^3]F&=S-6,\label{eq:coef3}\\
 [x^2]F&=e_2-3S+\sum_{i<j}D_{ij}+11,\label{eq:coef2}\\
 [x]F&=e_3-e_2+\sum_{\{i,j\}\in\pairs}D_{ij}(z_k+z_\ell)
       -\sum_{i<j}D_{ij}+2S-2\sum_iT_i-6,\label{eq:coef1}\\
 [1]F&=M-2\sum_i z_iT_i+6U.\label{eq:coef0}
\end{align}
These four equations will be solved exactly, while maintaining positivity
in Lemma~\ref{lem:venn}.

\section{A rational cancellation profile}\label{sec:construction}

Fix arbitrary $a,b,c,d\in\Z$. In this section $t$ is any positive
multiple of $52$; its sufficient size will be established later. Put
\begin{equation}\label{eq:Sdelta}
 S=a+6,\qquad \Delta=2-S=-a-4,
\end{equation}
and choose
\begin{equation}\label{eq:z}
 (z_1,z_2,z_3,z_4)=(t+1,\ t,\ -t,\ -t+S-1).
\end{equation}
In particular, $\sum z_i=S$ and $z_1-z_2=1$.

For an integer $r$, $\bit r$ denotes its least nonnegative residue modulo
$2$. Start with the following six bits:
\begin{center}
\begin{tabular}{c cccccc}
\toprule
 & $\varepsilon_{12}$ & $\varepsilon_{13}$ & $\varepsilon_{14}$
 & $\varepsilon_{23}$ & $\varepsilon_{24}$ & $\varepsilon_{34}$\\
\midrule
$a$ odd & $1$ & $\bit{1+c}$ & $0$ & $\bit{b+c+d}$ & $0$ & $\bit d$\\
$a$ even & $1$ & $\bit{b+d+1}$ & $0$ & $\bit{c+1}$ & $0$ & $\bit{c+d+1}$\\
\bottomrule
\end{tabular}
\end{center}
Let $K=b+2a+2$. Set $\delta_{ij}=\varepsilon_{ij}$ except at $24$, and set
\begin{equation}\label{eq:deltacorrection}
 \delta_{24}=K-\sum_{\{i,j\}\ne\{2,4\}}\varepsilon_{ij}.
\end{equation}
The six initial bits sum to $b$ modulo $2$, so $\delta_{24}$ is even.
Thus every $\delta_{ij}$ has the prescribed parity and
$\sum\delta_{ij}=K$.

Define
\begin{align}
 D_{12}&=\frac{2}{13}t^2-t+\delta_{12},
 &D_{34}&=\frac{2}{13}t^2+t+\delta_{34},\notag\\
 D_{13}&=\frac{6}{13}t^2+4\Delta t+\delta_{13},
 &D_{24}&=\frac{6}{13}t^2+\delta_{24},\label{eq:D}\\
 D_{14}&=\frac{5}{13}t^2-3\Delta t+\delta_{14},
 &D_{23}&=\frac{5}{13}t^2+\delta_{23}.\notag
\end{align}
All these quantities are integers. The individual corrections need not be
positive; positivity is required of the final atom counts, and will be
proved for those counts.

\begin{lemma}[The first two equations and the two parity conditions]
\label{lem:parity}
For \eqref{eq:z} and \eqref{eq:D}, the cubic and quadratic coefficients
in \eqref{eq:coef3}--\eqref{eq:coef2} are $a,b$. Moreover
\begin{equation}\label{eq:Bsum}
 B_*:=\frac12\left(e_3-e_2+
 \sum_{\{i,j\}\in\pairs}D_{ij}(z_k+z_\ell)
 -\sum_{i<j}D_{ij}+2S-6-c\right)
\end{equation}
is an integer, and $M\equiv d\pmod2$.
\end{lemma}
\begin{proof}
The cubic coefficient is $S-6=a$. Directly from \eqref{eq:z},
\[
 e_2=-2t^2-\Delta t+S-1,
 \qquad \sum D_{ij}=2t^2+\Delta t+K.
\]
Substitution in \eqref{eq:coef2} gives $K-2S+10=b$.

Since $t$ is a multiple of $52$, its contributions to all the $D_{ij}$
are even. Also $D_{24}$ is even and $D_{14}$ is even. If $a$ is odd,
$z\equiv(1,0,0,0)\pmod2$. The numerator in \eqref{eq:Bsum} reduces to
$D_{12}+D_{13}+D_{14}-c\equiv0$, while
$M\equiv D_{12}D_{34}\equiv d$.
If $a$ is even, $z\equiv(1,0,0,1)\pmod2$. The same numerator reduces
to $1+D_{14}+D_{23}-c\equiv0$, and
$M\equiv D_{23}+D_{12}D_{34}\equiv d$.
These are precisely the two rows of the parity table.
\end{proof}

\begin{lemma}[Cancellation]\label{lem:cancellation}
For fixed $f$, as $t$ tends to infinity through positive multiples of $52$,
\begin{equation}\label{eq:cancellation}
 B_*=2t^2+O_f(t),\qquad M=O_f(t^2).
\end{equation}
In particular, the $t^4$ and $t^3$ coefficients of $M$ vanish identically.
\end{lemma}
\begin{proof}
Here is a full expansion that also records the dependence on the input.
Write
\[
 u=\delta_{13},\quad v=\delta_{23},\quad w=\delta_{34},\quad
 r=\delta_{24}=b+2a+1-u-v-w.
\]
Then $u,v,w\in\{0,1\}$, and substitution in
\eqref{eq:Bsum} and \eqref{eq:M} gives
\begin{align}
 B_*&=2t^2+h_1t+h_0,\label{eq:Bexpansion}\\
 h_1&=-2a^2-\frac{37}{2}a+w-43,\notag\\
 h_0&=\frac{a(u+v)}2+a-\frac c2+2u+\frac52v+\frac52,\notag\\[2pt]
 M&=m_2t^2+m_1t+m_0,\label{eq:Mexpansion}\\
 m_2&=-4a^2-\frac{463}{13}a+
       \frac{-1000-7u-8v-7r+15w}{13},\notag\\
 m_1&=-a+au+4av-4ar-4+5u+16v-17r,\notag\\
 m_0&=av+w+ur+5v.\notag
\end{align}
These identities prove the lemma.

For clarity, the leading cancellation has a small rational origin.
Ignoring lower terms in $z$, take $z/t=(1,1,-1,-1)$ and take the
complementary-pair coefficients of $D/t^2$ to be $p,q_0,r_0$.
The quadratic-coefficient equation requires $p+q_0+r_0=1$.
The leading coefficient of $M$ then becomes
\[
 -1+4p+p^2+q_0^2+r_0^2.
\]
The rational choice
$(p,q_0,r_0)=(2,6,5)/13$ makes this zero. It is the rational point
$3(p+1)^2+(q_0-r_0)^2=4$ with
$(p+1,q_0-r_0)=(15,1)/13$.
The linear terms in \eqref{eq:D} cancel the next, cubic term of $M$
and give the coefficient $2$ in \eqref{eq:Bexpansion}.
\end{proof}

\section{Solving the last two equations over the integers}\label{sec:lift}

Set
\begin{equation}\label{eq:baseline}
 L=\frac7{13}t^2,\qquad T_0=\frac12t^2,\qquad
 U_0=\frac{27}{52}t^2,
\end{equation}
which are integers, and define
\begin{equation}\label{eq:W}
 W=\frac{M+6U_0-d}{2},\qquad h=B_*-2T_0.
\end{equation}
Lemma~\ref{lem:parity} makes both $W$ and $h$ integers. Let
\begin{equation}\label{eq:theta}
 \theta=\frac{W-(z_2+z_4)T_0-z_3h}{z_1-z_3}
       =\frac{W-(z_2+z_4)T_0+th}{2t+1}.
\end{equation}
First put
\begin{equation}\label{eq:Tpre}
 (\widetilde T_1,\widetilde T_2,\widetilde T_3,\widetilde T_4)
 =\bigl(\lfloor\theta\rfloor,\ T_0,\ h-\lfloor\theta\rfloor,\ T_0\bigr).
\end{equation}
Define
\begin{equation}\label{eq:R}
 R=W-\sum_i z_i\widetilde T_i,
 \qquad \eta=R-3\lfloor R/3\rfloor\in\{0,1,2\}.
\end{equation}
Finally set
\begin{align}
 T_1&=\widetilde T_1+\eta,&
 T_2&=\widetilde T_2-\eta,&
 T_3&=\widetilde T_3,&
 T_4&=\widetilde T_4,\label{eq:Tfinal}\\
 U&=U_0+\frac{\eta-R}{3}=U_0-\lfloor R/3\rfloor.\label{eq:Ufinal}
\end{align}

\begin{lemma}[Exact integer lift]\label{lem:lift}
All parameters in \eqref{eq:Tfinal}--\eqref{eq:Ufinal} are integers,
and they satisfy
\begin{equation}\label{eq:twolast}
 \sum_iT_i=B_*,\qquad 2\sum_i z_iT_i-6U=M-d.
\end{equation}
Moreover $0\le R\le2t$ and $|U-U_0|\le t$.
\end{lemma}
\begin{proof}
The definition of $\theta$ is the exact real solution of
$T_1+T_3=h$ and
$z_1T_1+z_3T_3=W-(z_2+z_4)T_0$.
Thus
\[
 R=(2t+1)(\theta-\lfloor\theta\rfloor)
\]
is an integer in $\{0,1,\ldots,2t\}$.
The preliminary $\widetilde T_i$ sum to $B_*$; the correction preserves
that sum. Because $z_1-z_2=1$, it increases their weighted sum by
exactly $\eta$. Consequently
\[
 \sum_i z_iT_i=W-R+\eta=W+3(U-U_0),
\]
which is the second equation in \eqref{eq:twolast}.
The definition of $\eta$ ensures integrality of $U$, and the residual
bound gives $|U-U_0|\le2t/3\le t$.
\end{proof}

Together, Lemmas~\ref{lem:parity} and~\ref{lem:lift} imply the exact
polynomial identity
\begin{equation}\label{eq:Fisf}
 F(x)=x^4+ax^3+bx^2+cx+d=f(x).
\end{equation}
It remains essential to prove that the parameters define an actual graph.

\section{Every palette atom has a positive margin}\label{sec:positive}

\begin{lemma}[Uniform leading profile]\label{lem:positive}
For each fixed $f$, the atom counts in
\eqref{eq:atom4}--\eqref{eq:atom1}, with the parameters constructed above,
satisfy
\begin{equation}\label{eq:atomprofile}
 m_J=\lambda_Jt^2+O_f(t),
\end{equation}
where all the constants $\lambda_J$ are strictly positive:
\begin{center}
\begin{tabular}{lc}
\toprule
Atom type & $\lambda_J$\\
\midrule
Each singleton $\{i\}$ & $1/52$\\
Each triple $[4]\setminus\{i\}$ & $1/52$\\
$[4]$ & $1/52$\\
$\{1,2\},\{3,4\}$ & $17/52$\\
$\{1,3\},\{2,4\}$ & $1/52$\\
$\{1,4\},\{2,3\}$ & $5/52$\\
\bottomrule
\end{tabular}
\end{center}
Furthermore
\begin{equation}\label{eq:Nnasymptotic}
 N=\frac{55}{52}t^2+O_f(t),\qquad
 n:=N-L=\frac{27}{52}t^2+O_f(t).
\end{equation}
\end{lemma}
\begin{proof}
Write $E=B_*-4T_0$. From Lemma~\ref{lem:cancellation}, $E=O_f(t)$
and $W-ST_0=O_f(t^2)$. Equation~\eqref{eq:theta} can be rewritten as
\begin{equation}\label{eq:thetadiff}
 \theta-T_0=\frac{W-ST_0+tE}{2t+1}=O_f(t).
\end{equation}
Rounding and the bounded correction $\eta$ therefore give
$T_i=T_0+O_f(t)$ for every $i$. Lemma~\ref{lem:lift} gives
$U=U_0+O(t)$.

The quadruple and triple leading coefficients are respectively
$7/13-27/52=1/52$ and $27/52-1/2=1/52$.
For a pair with $D_{ij}/t^2\longrightarrow\beta_{ij}$, the coefficient
is $1-\beta_{ij}-27/52$, giving $17/52,1/52,5/52$ in the three
cases of \eqref{eq:D}. At a singleton, the sum of the three incident
$D_{ij}$ has leading coefficient $(2+6+5)/13=1$. The singleton
coefficient is therefore $1-3/2+27/52=1/52$.

Summing the fifteen coefficients gives $55/52$. Alternatively,
\eqref{eq:N} and $\sum D_{ij}=2t^2+O_f(t)$,
$\sum T_i=2t^2+O_f(t)$ immediately yield \eqref{eq:Nnasymptotic}.
\end{proof}

\begin{proof}[Proof of the existence assertion in Theorem~\ref{thm:main}]
Choose a sufficiently large positive multiple $t$ of $52$.
Lemma~\ref{lem:positive} makes all fifteen integer atom counts positive,
and also makes $n=N-L$ a positive integer. Lemma~\ref{lem:venn}
constructs actual subsets $A_i$ of an $N$-element set. Form the two-clique
graph described in Section~\ref{sec:palettes}. Equations
\eqref{eq:Pfactor} and \eqref{eq:Fisf} give
$P_G(q)=\fall qNf(q-n)$.

The four positive singleton atoms supply four distinct old colors for
$v_1,\ldots,v_4$, so the graph is $N$-colorable. Since it contains
$K_N$, its chromatic number is $N$. It is connected: for example, a
vertex of type $\{1\}$ is adjacent to $v_2$, joining the two cliques.
The explicit choice of $t$ and the bounds in \eqref{eq:size} are proved
next.
\end{proof}

\section{An explicit size bound}\label{sec:effective}

The preceding argument already closes all existence quantifiers. This
section supplies a predetermined input-dependent value of $t$, so even
the search for a sufficiently large scale can be omitted.

\begin{proposition}\label{prop:effective}
Let $\htf$ and $t$ be as in \eqref{eq:height}. Then every atom is strictly
positive, $0<n<t^2$, and $N+4\le2t^2$.
\end{proposition}
\begin{proof}
We give explicit bounds, keeping generous constants. Since $\htf\ge10$,
$u,v,w\in\{0,1\}$, and $r=b+2a+1-u-v-w$, we have
\[
 |S|\le\htf,\quad |\Delta|\le\htf,\quad |r|\le2\htf.
\]
The exact expansions \eqref{eq:Bexpansion}--\eqref{eq:Mexpansion} imply
\begin{equation}\label{eq:boundscoeff}
 |h_1|\le5\htf^2,\quad |h_0|\le4\htf,\quad
 |m_2|\le9\htf^2,\quad |m_1|\le21\htf^2,\quad |m_0|\le\htf^2.
\end{equation}
For example, $|h_1|\le2\htf^2+\tfrac{37}{2}\htf+44<5\htf^2$;
the remaining bounds follow directly by the triangle inequality in the
displayed formulas, using $\htf\ge10$.
For $t\ge1$, these give
\[
 |E|=|B_*-2t^2|\le6\htf^2t,\qquad |M|\le31\htf^2t^2.
\]
As $W=(M+6U_0-d)/2$ and $T_0=t^2/2$, it follows that
\begin{equation}\label{eq:Wbound}
 |W-ST_0|\le20\htf^2t^2.
\end{equation}
Equation~\eqref{eq:thetadiff} consequently gives
$|\theta-T_0|\le13\htf^2t$. The exact real value of the third
coordinate is $T_0+E-(\theta-T_0)$, at distance at most
$19\htf^2t$ from $T_0$. Applying the floor and the correction
$\eta\in\{0,1,2\}$ gives the convenient common bound
\begin{equation}\label{eq:Tbound}
 |T_i-T_0|\le20\htf^2t\quad(1\le i\le4),\qquad |U-U_0|\le t.
\end{equation}

Our chosen $t$ is at least $\htf$ and at least $52$. Thus
$|z_i|\le2t$, and each $D_{ij}$ differs from its quadratic leading
term by at most $5\htf t$. Equations
\eqref{eq:atom4}--\eqref{eq:atom1} now bound the deviation from the
profile in Lemma~\ref{lem:positive} as follows:
\begin{center}
\begin{tabular}{lc}
\toprule
Atom type & An upper bound for $|m_J-\lambda_Jt^2|$\\
\midrule
Quadruple & $t$\\
Triple & $21\htf^2t$\\
Pair & $41\htf^2t$\\
Singleton & $62\htf^2t$\\
\bottomrule
\end{tabular}
\end{center}
For instance, the singleton error is at most
$2t+15\htf t+60\htf^2t+t<62\htf^2t$.
The common upper bound $70\htf^2t$ is therefore valid for every atom.
Since $t=5200\htf^2$ and every $\lambda_J\ge1/52$,
\begin{equation}\label{eq:strictmargin}
 m_J\ge\frac{t^2}{52}-70\htf^2t
      =30\htf^2t>0.
\end{equation}

Finally, \eqref{eq:N} gives the exact relation
\[
 n-U_0=S+\Delta t+K-E+(U-U_0).
\]
Using $|K|\le2\htf$ and the previous bounds yields
\begin{equation}\label{eq:nbound}
 \left|n-\frac{27}{52}t^2\right|\le7\htf^2t.
\end{equation}
Because $7\htf^2/t=7/5200$, this proves $0<n<t^2$.
Also
\[
 N+4\le\left(\frac{55}{52}+\frac7{5200}\right)t^2+4<2t^2.
\]
This completes the proof of Proposition~\ref{prop:effective} and
Theorem~\ref{thm:main}.
\end{proof}

\section{Arithmetic consequences and algorithmic content}

\begin{corollary}[Fields and splitting fields]\label{cor:fields}
Every quartic number field is generated by a chromatic root. Moreover,
the splitting field of every monic integer quartic is exactly the
splitting field of a chromatic polynomial of a connected graph.
\end{corollary}
\begin{proof}
Every number field has an integral primitive element. If
$K=\Q(\alpha)$ has degree four, apply Theorem~\ref{thm:main} to its
monic minimal polynomial. Then $\alpha+n$ is a chromatic root and
$\Q(\alpha+n)=\Q(\alpha)=K$.
For the second assertion, translation by $n\in\Z$ preserves a
polynomial's splitting field, and the extra factor $\fall qN$ splits
over $\Q$. Equation~\eqref{eq:main} therefore gives equality of
splitting fields, and hence the same Galois group over $\Q$.
\end{proof}

The second assertion proves the degree-four instance of the splitting-field
realization statement posed in \cite[Conjecture 7.1]{Gilmore}.
Indeed, after making a rational quartic $g$ monic, choose a positive
integer $m$ clearing its coefficient denominators. The polynomial
$m^4g(x/m)$ is monic and integral and has the same splitting field as $g$.
The distinction between a root field and a splitting field matters:
a nonnormal quartic field is realized as a root field, not asserted to
be itself a splitting field.

\begin{corollary}[A bounded arithmetic construction]\label{cor:algorithm}
There is an algorithm that, on input $(a,b,c,d)\in\Z^4$, produces a
positive shift and a fifteen-integer description of $G$ using a bounded
number of exact arithmetic operations. The integers in this compressed
description have $O(\log\htf)$ bits.
\end{corollary}
\begin{proof}
Use the predetermined $t=5200\htf^2$ and the displayed formulas, including
the two integer divisions in \eqref{eq:Tpre} and \eqref{eq:R}.
All intermediate quantities are polynomials of bounded degree in the
input and $t$, followed by division by positive explicitly known integers.
Their bit lengths are therefore $O(\log\htf)$. The graph is specified
by its two cliques and the fifteen atom multiplicities.
\end{proof}

The stated operation count concerns the compressed description. Writing
out all vertices and edges has the corresponding output cost.
For small examples one may try $t=52,104,208,\ldots$ and stop when the
fifteen atoms and $n$ are positive. Theorem~\ref{thm:main} proves eventual
termination, but no such search is required by the theorem.

\section{An exact example and independent verification}\label{sec:example}

Take $f(x)=x^4+x+1$. The construction already works at $t=3328$, which
is much smaller than the uniform sufficient value. It gives
\[
 L=5\,963\,776,\qquad N=11\,840\,301,\qquad n=5\,876\,525.
\]
The complete graph certificate consists of the following atom counts.
An element of type $J$ is nonadjacent to precisely the small-clique
vertices indexed by $J$.
\begin{center}
\begin{tabular}{cr@{\qquad}cr@{\qquad}cr}
\toprule
$J$ & $m_J$ & $J$ & $m_J$ & $J$ & $m_J$\\
\midrule
$1$ & 203\,420 & $12$ & 3\,620\,454 & $123$ & 212\,263\\
$2$ & 352\,037 & $13$ & 266\,968 & $124$ & 216\,729\\
$3$ & 294\,324 & $14$ & 1\,021\,286 & $134$ & 212\,264\\
$4$ & 391\,979 & $23$ & 930\,382 & $234$ & 347\,570\\
 & & $24$ & 73\,948 & $1234$ & 213\,721\\
 & & $34$ & 3\,482\,956 & &\\
\bottomrule
\end{tabular}
\end{center}
Every entry is positive, and the entries sum to $N$. This defines a
specific graph on $11\,840\,305$ vertices and proves the exact identity
\[
 P_G(q)=\fall q{11\,840\,301}
 \left((q-5\,876\,525)^4+(q-5\,876\,525)+1\right).
\]
No minimality claim is made for its order or shift.

There is a counting check that does not use the intersection-coefficient
equations of Section~\ref{sec:palettes}. Let $B=\overline G$, a bipartite
graph with parts of sizes $4,N$, and let $r_k$ count its matchings of
size $k$. A coloring of $G$ can identify at most one vertex of each
clique in a color class, so the standard matching expansion is
\begin{equation}\label{eq:rook}
 P_G(q)=\sum_{k=0}^4r_k\fall q{N+4-k}
       =\fall qN\sum_{k=0}^4r_k\fall{q-N}{4-k}.
\end{equation}
This is also the biclique matching formula of \cite[Section 2.1]{Bohn}.
For the example, direct injection counting on the fifteen atom types gives
\begin{align*}
 r_0&=1,\\
 r_1&=23\,855\,110,\\
 r_2&=213\,399\,816\,634\,375,\\
 r_3&=848\,445\,732\,099\,071\,934\,466,\\
 r_4&=1\,264\,984\,755\,431\,613\,627\,400\,650\,753.
\end{align*}
Substitution in \eqref{eq:rook} gives $f(y+L)$ exactly, with $y=q-N$.

The implementation in Appendix~\ref{app:code} was checked on the displayed
example and four additional named inputs, on 1000 deterministic-seed
random inputs with coefficient sizes up to $10^8$, and by a separately
written injection-counting routine on 256 further inputs with coefficient
sizes up to $10^{18}$. Separate symbolic expansions checked the cancellations
and all sixteen coefficient-parity patterns. These are implementation
checks; the universal proof is the explicit construction and inequalities
above, and uses no probabilistic inference.

\section{Scope, provenance, and priority}\label{sec:scope}

The result is the unrestricted \emph{quartic case}: there is no restriction
on coefficient signs, coefficient height, real signature, Galois group,
or the particular quartic algebraic integer. The same construction also
handles reducible monic quartics. Nothing in the proof resolves the
degree-five or arbitrary-degree case of the $\alpha+n$ conjecture.

The original problem and the biclique framework are due to the cited
literature. The new mathematical ingredient presented here is the
specific rational cancellation profile and its positive integral lift,
including a uniform quartic height bound. The supplied corpus was inspected
at commit \texttt{9f8547bfd3b8f6560f20c1b63028dc87a07a2e33}; the
relevant proof-design connection is manuscript~245, not an assumed theorem
about chromatic roots.

The literature assessment consulted the full texts of
\cite{Bohn,CM,Gilmore,DM22}, the official abstract of \cite{DM19}, and
Cameron's conjecture list dated 2 November 2026 \cite{CameronWeb}.
The full text of \cite{DM19} was not available during preparation; it has
therefore not been checked line by line for overlap of intermediate
constructions. No identical general quartic solution was identified in
the consulted material. This is not a certification of absolute publication
priority, and the manuscript has not undergone external peer review.
The algebraic proof is provided in full so that correctness can be assessed
independently of those bibliographic questions.

\appendix
\section{Self-contained exact-arithmetic implementation}\label{app:code}

The following Python 3 code uses only the standard library. Its input
is four integers $(a,b,c,d)$. The deterministic function
\texttt{construct\_bounded} uses the proved bound; \texttt{construct}
tries smaller powers-of-two scales for convenience. The independent
\texttt{rook\_numbers} routine counts matchings from the atom multiplicities.
It does not call the centered coefficient formula. An integer mask
$1\le m\le15$ denotes $\{i+1: \text{bit }i\text{ of }m\text{ is }1\}$.
For example, mask $10$ denotes $\{2,4\}$.

\begin{lstlisting}
"""Exact quartic biclique construction. Research verification, 2026-11-28.

All graph counts are arbitrary-precision integers.  A column of type mask
means a vertex in a large clique whose NON-neighbours in the four-clique
are the rows in mask.  No expansion of the large graph is required.
"""
from itertools import combinations
from math import prod
import json

PAIRS = list(combinations(range(4), 2))

def elementary(z, k):
    return sum(prod(z[i] for i in I) for I in combinations(range(4), k))

def make_certificate(a, b, c, d, t):
    assert t > 0 and t % 52 == 0
    S = a + 6
    delta = 2 - S
    z = [t + 1, t, -t, -t + S - 1]
    if a % 2:
        eps = {(0,1):1, (0,2):(1+c)%2, (0,3):0,
               (1,2):(b+c+d)%2, (1,3):0, (2,3):d%2}
    else:
        eps = {(0,1):1, (0,2):(b+d+1)%2, (0,3):0,
               (1,2):(c+1)%2, (1,3):0, (2,3):(d+c+1)%2}
    K = b + 2*a + 2
    correction = K - sum(eps.values())
    assert correction % 2 == 0
    eps[(1,3)] += correction
    D = {
        (0,1):2*t*t//13 - t + eps[(0,1)],
        (0,2):6*t*t//13 + 4*delta*t + eps[(0,2)],
        (0,3):5*t*t//13 - 3*delta*t + eps[(0,3)],
        (1,2):5*t*t//13 + eps[(1,2)],
        (1,3):6*t*t//13 + eps[(1,3)],
        (2,3):2*t*t//13 + t + eps[(2,3)],
    }
    def comp(I): return [i for i in range(4) if i not in I]
    e2, e3, e4 = [elementary(z,k) for k in (2,3,4)]
    Hnum = e3-e2+sum(D[I]*sum(z[j] for j in comp(I)) for I in PAIRS)-sum(D.values())+2*S-6-c
    assert Hnum % 2 == 0
    H = Hnum//2
    M = e4 + sum(D[I]*prod(z[j] for j in comp(I)) for I in PAIRS)
    M += D[(0,1)]*D[(2,3)] + D[(0,2)]*D[(1,3)] + D[(0,3)]*D[(1,2)]
    assert (M-d) % 2 == 0
    L = 7*t*t//13
    T0 = t*t//2
    U0 = 27*t*t//52
    W = (M+6*U0-d)//2
    h = H - 2*T0
    numerator = W-(z[1]+z[3])*T0-z[2]*h
    denominator = z[0]-z[2]
    T1 = numerator//denominator
    T = [T1,T0,h-T1,T0]
    R = W-sum(zi*Ti for zi,Ti in zip(z,T))
    eta = R % 3
    T[0] += eta
    T[1] -= eta
    U = U0+(eta-R)//3
    assert sum(T) == H
    assert 2*sum(zi*Ti for zi,Ti in zip(z,T))-6*U == M-d
    atoms = {15:L-U}
    for i in range(4):
        atoms[15 ^ (1<<i)] = U-T[i]
    for i,j in PAIRS:
        k,l = comp((i,j))
        atoms[(1<<i)|(1<<j)] = T[k]+T[l]-D[(i,j)]-U
    for i in range(4):
        atoms[1<<i] = z[i]-sum(T[j] for j in range(4) if j != i)+sum(v for I,v in D.items() if i in I)+U
    N = sum(atoms.values())
    n = N-L
    assert N == L+S+sum(D.values())-sum(T)+U
    centered = [1,S-6,e2-3*S+sum(D.values())+11,
                e3-e2+sum(D[I]*sum(z[j] for j in comp(I)) for I in PAIRS)-sum(D.values())+2*S-2*sum(T)-6,
                M-2*sum(zi*Ti for zi,Ti in zip(z,T))+6*U]
    assert centered == [1,a,b,c,d], (centered,[1,a,b,c,d])
    return dict(coefficients=[1,a,b,c,d],t=t,z=z,D={str(I):v for I,v in D.items()},T=T,U=U,L=L,N=N,n=n,
                atoms={str(m):atoms[m] for m in range(1,16)},minimum_atom=min(atoms.values()),
                valid=min(atoms.values())>0 and n>0)

def rook_numbers(atomdict):
    """Independent injection DP on the four labelled rows."""
    dp=[0]*16
    dp[0]=1
    for maskstr,count in atomdict.items():
        mask=int(maskstr)
        assert count>=0
        ff=[1]
        for k in range(1,5): ff.append(ff[-1]*(count-k+1))
        ndp=[0]*16
        for used,ways in enumerate(dp):
            if not ways: continue
            available=mask & (15^used)
            chosen=available
            while True:
                ndp[used|chosen]+=ways*ff[chosen.bit_count()]
                if chosen==0: break
                chosen=(chosen-1)&available
        dp=ndp
    return [sum(dp[mask] for mask in range(16) if mask.bit_count()==k) for k in range(5)]

def verify_independently(cert):
    assert cert['valid']
    m=rook_numbers(cert['atoms'])
    got=[1,m[1]-6,m[2]-3*m[1]+11,m[3]-m[2]+2*m[1]-6,m[4]]
    _,a,b,c,d=cert['coefficients']
    L=cert['L']
    want=[1,4*L+a,6*L*L+3*a*L+b,4*L**3+3*a*L*L+2*b*L+c,L**4+a*L**3+b*L*L+c*L+d]
    assert got==want, (got,want)
    return m

def construct(a,b,c,d):
    t=52
    while True:
        cert=make_certificate(a,b,c,d,t)
        if cert['valid']:
            cert['rook_numbers']=verify_independently(cert)
            return cert
        t*=2

def construct_bounded(a,b,c,d):
    height=10+abs(a)+abs(b)+abs(c)+abs(d)
    cert=make_certificate(a,b,c,d,5200*height*height)
    assert cert['valid']
    assert cert['minimum_atom'] > 0
    assert 0 < cert['n'] < cert['t']**2
    assert cert['N']+4 <= 2*cert['t']**2
    cert['rook_numbers']=verify_independently(cert)
    return cert

if __name__ == '__main__':
    # A smaller exact certificate for x^4+x+1:
    cert=make_certificate(0,0,1,1,3328)
    cert['rook_numbers']=verify_independently(cert)
    print(json.dumps(cert,sort_keys=True,indent=2))
    # The uniform construction requires no search:
    bounded=construct_bounded(0,0,1,1)
    print('Bounded construction verified:',bounded['N']+4)
\end{lstlisting}

\begin{thebibliography}{9}

\bibitem{CM}
P. J. Cameron and K. Morgan,
\emph{Algebraic properties of chromatic roots},
Electronic Journal of Combinatorics \textbf{24}(1) (2017), P1.21.
\href{https://doi.org/10.37236/6578}{doi:10.37236/6578}.
\href{https://emis.de/ft/29131}{Full text}.

\bibitem{Bohn}
A. Bohn, \emph{Chromatic roots as algebraic integers},
Discrete Mathematics and Theoretical Computer Science Proceedings
\textbf{AR} (2012), 543--554.
\href{https://doi.org/10.46298/dmtcs.3061}{doi:10.46298/dmtcs.3061}.
\href{https://www.math.nagoya-u.ac.jp/fpsac12/download/contributed/dmAR0148.pdf}{Full text}.

\bibitem{Gilmore}
H. J. Gilmore,
\emph{Algebraic Properties of Chromatic Polynomials and Their Roots},
MSc thesis, University of Waikato, 2015.
\url{https://hdl.handle.net/10289/9367}.

\bibitem{DM19}
D. Delbourgo and K. Morgan,
\emph{An algorithm which outputs a graph with a specified chromatic factor},
Discrete Applied Mathematics \textbf{257} (2019), 128--150.
\href{https://doi.org/10.1016/j.dam.2018.10.033}{doi:10.1016/j.dam.2018.10.033}.
Official abstract consulted through the author's
\href{https://profiles.auckland.ac.nz/daniel-delbourgo/publications}{University of Auckland publication record}.

\bibitem{DM22}
D. Delbourgo and K. Morgan, \emph{On the fifth chromatic coefficient},
Australasian Journal of Combinatorics \textbf{84}(1) (2022), 56--85.
\url{https://ajc.maths.uq.edu.au/pdf/84/ajc_v84_p056.pdf}.

\bibitem{VMRC}
Verifiable Mathematical Reasoning Corpus,
manuscript 245, \emph{Roots and coefficient peaks of graph polynomials},
research manuscript dated 21 November 2026, Sections 2--3.
Snapshot \texttt{9f8547bfd3b8f6560f20c1b63028dc87a07a2e33},
accessed 28 November 2026.
\url{https://github.com/QEEDDD/Verifiable-Mathematical-Reasoning-Corpus}.

\bibitem{CameronWeb}
P. J. Cameron, \emph{Peter Cameron's conjectures},
page dated 2 November 2026, accessed 28 November 2026.
\url{https://cameroncounts.github.io/web/problems/cameronconjs.html}.

\end{thebibliography}
\end{document}
